% ============================================================================
% CHAPTER 3 — METHODOLOGY
% Source to adapt: thesisproposal/main.tex §4 (Proposed Methodology) — the
% two-phase architecture is implemented as described; this chapter documents
% the system AS BUILT (src/tdbf/agent/*, scripts/07_deployable_pipeline.py).
% Figures available: ../figures/tdbf_pipeline.pdf, ../figures/tdbf_agents.pdf
% ============================================================================
\chapter{Methodology}
\label{ch:methodology}

\section{Overview: the Test-First Pipeline}
\label{sec:method-overview}
% TODO: adapt proposal §4.1. The strict order as implemented: BRT agent →
% verify the test FAILS on the buggy revision (gold-free gate) → fix agent →
% offline scoring. Include figure: \includegraphics{tdbf_pipeline.pdf}

\section{The Gold-Free Invariant}
\label{sec:method-goldfree}
% TODO: NEW section (the design contribution that emerged during the work).
% No agent ever holds a tool that observes the gold patch, the gold-fixed
% tree, or the held-out FAIL_TO_PASS tests; gold is an offline measurement
% instrument only. Enforced programmatically: allowlisted tool registries +
% assert_gold_free / assert_fix_gold_free guards; context_safe_dict() as the
% only instance view an agent receives. Why it matters: a pipeline whose
% decisions depend on the gold patch cannot deploy — this is the difference
% between an oracle study and an honest system.

\section{Agent Runtime and Tools}
\label{sec:method-runtime}
% TODO: describe the vendor-agnostic tool-use loop (runtime.py): terminal
% submit_* tools, iteration/token caps, persistent per-instance container
% (session.py). Tool tables:
%   BRT agent:  read_file, grep, glob, run_python, run_test, submit_test
%   Fix agent:  read_file, grep, glob, edit_file, run_brt, run_tests,
%               revert, submit_fix   (run_brt has the BRT baked in — the
%               fix agent cannot alter its own success criterion)

\section{Phase 1: BRT Generation and Validation}
\label{sec:method-phase1}
% TODO: adapt proposal §4.2.1 to as-built: the agent probes the buggy tree,
% writes the test, runs it; the gate accepts iff the test reproduces a clean
% failure on buggy (candidate = L3∪L4 — the split is precisely what cannot
% be told without gold, so the deployable pipeline must fix L3s too).

\section{Phase 2: Fix Generation}
\label{sec:method-phase2}
% TODO: adapt proposal §4.2.2 to as-built: BRT as a GUIDE not the oracle
% (prompt design), bail-out rule, targeted regression testing, revert
% discipline, fix patch = git diff of tracked tree (BRT auto-excluded).

\section{The Reactive Baseline (Arm A)}
\label{sec:method-baseline}
% TODO: same model, tools, budgets, container, scoring; run_brt removed,
% run_tests kept; no Phase 1. Quote proposal fairness rule: the baseline
% receives only issue text, repository context, and repository tests.

\section{BRT Classification (L0--L4) and Metrics}
\label{sec:method-levels}
% TODO: the L0–L4 taxonomy and its SWE-bench-vocabulary mapping
% (L4 = Fail→Pass = plausible BRT; L3 = Fail→Fail = candidate-not-plausible);
% C2P = L4/(L3+L4); resolved = held-out FAIL_TO_PASS all pass ∧ PASS_TO_PASS
% preserved. Table format from webapp seed doc "What TDBF is".
